\chapter{Два дрона}\label{ch:dva-drona}

\textbf{Что вы узнете:} как написать шоу для двух дронов, не столкнув их.

\section{Два блока}

Для второго дрона просто добавьте ещё один блок \kod{d = show.drone(2)} и его команды. Каждый дрон выполняет свою программу, они не знают друг о друге. Все дроны стартуют одновременно.

\solfile{l10_pass.py}

Дроны стоят в клетках 13 и 16 и меняются местами. Они летят по одной линии навстречу друг другу, поэтому один летит на 80\cm, а другой на 150\cm.

\begin{figure}[H]
  \centering
  % Вид сбоку: два дрона разминуются на разной высоте
\begin{tikzpicture}[x=0.075cm, y=0.0275cm, font=\sffamily\small]
  \draw[-{Stealth}, ink] (-5,0) -- (150,0) node[right] {$x$, см};
  \draw[-{Stealth}, ink] (0,0) -- (0,172) node[above] {$z$, см};
  \draw[dashed, rule] (-5,150) -- (140,150) node[right, font=\scriptsize] {лимит 150};
  \draw[ledblue, very thick, -{Stealth}] (0,80) -- (120,80);
  \draw[ledred, very thick, -{Stealth}] (120,150) -- (0,150);
  \node[ledblue, anchor=north] at (30,76) {дрон 1, 80\,см};
  \node[ledred, anchor=south] at (30,154) {дрон 2, 150\,см};
  \draw[<->, accent, thick] (60,82) -- (60,148) node[midway, right, text=accent, font=\scriptsize\bfseries] {70\,см};
\end{tikzpicture}

  \caption{Вид сбоку: дроны проходят друг над другом с зазором 70\cm.}\label{fig:heights}
\end{figure}

\section{Правила безопасной дистанции}

Расстояние между дронами считается в трёх измерениях, вместе с высотой. Проверка требует, чтобы оно \textbf{всегда было не меньше 60\cm}.

\begin{itemize}
  \item \textbf{Старт.} Дроны стоят на земле: стартовые клетки соседних дронов держите не ближе двух клеток (80\cm).
  \item \textbf{Полёт.} Если пути пересекаются, разведите дронов по высоте не меньше чем на 60\cm; лучше 70 и больше.
  \item \textbf{Посадка.} Конечные клетки тоже не ближе двух клеток.
\end{itemize}

Вот что будет, если оба дрона летят на одной высоте:

\solfile{bad_examples/bad_close.py}

Проверка скажет: «Дроны 1 и 2 сближаются до 2 см на 12,3 с (порог 60 см)». Можно передвинуть время на 12,3\,с и посмотреть на карте, как дроны проходят друг через друга.

\begin{zametka}
Чтобы развести дронов, есть три способа: разные высоты, разное время (один ждёт, пока пролетит другой) и разные маршруты. Первые два приёма в этой книге: высоты здесь, ожидание в уроке~\ref{ch:ochered}.
\end{zametka}

\begin{probuy}
  \item Два дрона летят рядом, не пересекаясь: 13 $\to$ 49 и 16 $\to$ 52, оба на 100\cm. Почему здесь не нужны разные высоты?
  \item «Крест»: дрон 1 летит из клетки 13 в 52, дрон 2 из 16 в 49. Они пересекаются в центре. Разведите их по высоте (80 и 150\cm).
  \item Дрон 1 облетает квадрат (клетки 13, 16, 52, 49), а дрон 2 висит выше в клетке 27 и мигает красным четыре раза. Подберите высоты, чтобы проверка прошла.
  \item Исправьте файл \texttt{bad\_close.py} выше.
\end{probuy}
\otvety{10}
